\item Una placa de hierro se mantiene contra una rueda de hierro tal que una fuerza de fricción cinética de 50.0 N actúa entre los dos pedazos de metal. La rapidez relativa a la cual deslizan entre sí es de 40.0 m/s. (a) Calcule la rapidez con la que se convierte energía mecánica a energía interna. (b) La placa y la rueda tienen, cada uno, una masa de 5.00 kg, y cada uno recibe 50\% de la energía interna. Si el sistema opera, como se ha descrito, durante 10.0 s y a cada objeto le es permitido lograr una temperatura interna uniforme, ¿cuál es el incremento de temperatura resultante?
